\section{Data Model}
\label{sec:data_model}
This section defines the data model that we introduce as the basis of our mapping algebra. This data model is a version of the relational data model,
	restricted to a special kind of relations, that we call \emph{mapping relations}, in which the only possible values
	%restricted to values that
are RDF terms (plus a special value to capture processing errors).

As a preliminary step for defining the model formally, we introduce relevant concepts of RDF:
We let $\symAllStrings$ be the countably infinite set of strings (sequences of
Unicode code points) and $\symAllIRIs$ be the subset of $\symAllStrings$ that
consists of all strings that are valid IRIs. 
Moreover, $\symAllLiterals$ is a countably infinite set of pairs
$ \literalTuple \in \symAllStrings \times \symAllIRIs$, which we
call \emph{literals}. 
For each such literal~$\literal = \literalTuple$, $\lex$ is
its \emph{lexical form} and $\dt$ is its \emph{datatype
IRI}.\footnote{To avoid overly complex formulas in this paper, we ignore the
option of literals that have a language tag. Yet, the formalism can easily be
extended to cover language-tagged literals as well.}
Furthermore, we assume a countably infinite set~$\symAllBNodes$ of \emph{blank
nodes}, which is disjoint from both $\symAllStrings$ and $\symAllLiterals$. 
The set~$\symTermUniverse$ of all \emph{RDF terms} is defined as
$\symTermUniverse = \symAllIRIs \cup \symAllBNodes \cup \symAllLiterals$, and
we write $\symTermSeqUniverse$ to denote the set of all possible sequences of
terms; i.e., each element of $\symTermSeqUniverse$ is a sequence of elements of
$\symTermUniverse$.
As usual, an \emph{RDF triple} is a tuple $(s,p,o) \in (\symAllIRIs \cup \symAllBNodes) \times \symAllIRIs \times \symTermUniverse$, an \emph{RDF graph} is a set of RDF triples, and an \emph{RDF dataset} is a set $\big\{ G_\mathrm{dflt}, (n_1,G_1), \ldots, (n_m,G_m) \big\}$ where $G_\mathrm{dflt}, G_1, \ldots, G_m$ are RDF graphs, $m \geq 0$, and $n_i \in \symAllIRIs \cup \symAllBNodes$ for all $i \in \{1,\ldots,m\}$.

As additional ingredients for defining mapping relations, we assume a countably infinite set~$\symAttrUniverse$ of
\emph{attributes} and a special value $\error$ such that $\error \notin \symTermUniverse$%
	%, which shall be used to denote errors%
, and we define the notion of a \emph{mapping tuple}, as
	%the types of tuples
	shall be
used in mapping relations.

\begin{definition}
	\normalfont
	A \textbf{mapping tuple} is a partial function $\mappingTuple:
		\symAttrUniverse \rightarrow \symTermUniverse \cup \{\error\}$.
\end{definition}

\todo[inline]{Olaf, continue here ...}

Two mapping tuples $\mappingTuple$ and $\mappingTuple'$ are \emph{compatible} if, for every attribute~$\attr \in \fctDom{\mappingTuple} \cap \fctDom{\mappingTuple'}$, it holds that $\mappingTuple(\attr) = \mappingTuple'(\attr)$.
Given two mapping tuples $\mappingTuple$ and $\mappingTuple'$ that are compatible, we write $\mappingTuple \cup \mappingTuple'$ to denote the mapping tuple $\mappingTuple_\text{merged}$ that is the merge of $\mappingTuple$ and $\mappingTuple'$; that is, $\fctDom{\mappingTuple_\text{merged}} = \fctDom{\mappingTuple} \cup \fctDom{\mappingTuple'}$ and, for every attribute~$\attr \in \fctDom{\mappingTuple_\text{merged}}$,
\begin{equation*}
	\mappingTuple_\text{merged}(\attr) =
	\begin{cases}
		\mappingTuple(\attr) & \text{if $\attr \in \fctDom{\mappingTuple}$,} \\
		\mappingTuple'(\attr) & \text{else.}
	\end{cases}
\end{equation*}

Finally, our notion of a mapping relation is then defined as follows.

\begin{definition}
	\normalfont
	A \textbf{mapping relation}~$r$ is a tuple $\mappingRel$, where
	$\symAttrSubSet \subset \symAttrUniverse$ is a finite, non-empty set of attributes
    and 	
	$\symMappingInst$ is a set of mapping tuples such
		that, for every such tuple~$\mappingTuple \in \symMappingInst$, it holds that $\mathrm{dom}(\mappingTuple) = \symAttrSubSet$.
		%that $\mathrm{dom}(\mappingTuple) = \symAttrSubSet$ for all $\mappingTuple$ in $\symMappingInst$.
	We call $\symAttrSubSet$ the \textbf{schema} of the mapping relation
	and $\symMappingInst$ is the \textbf{instance} of the mapping relation.
	\label{def:mapping_relation}
\end{definition}

For the sake of conciseness, we sometimes write $\mappingTuple \in r$ to denote that $\mappingTuple$ is a mapping tuple in
	the instance~%
$\symMappingInst$ of a mapping relation $r=\mappingRel$.

\todo[inline]{TODO (Olaf): here comes the part about the RDF datasets generated from mapping relations}


we assume four designated attributes: $\subjAttr \in \symAttrUniverse$,
$\predAttr \in \symAttrUniverse$, $\objAttr \in
\symAttrUniverse$, and $\graphAttr \in \symAttrUniverse$
